\documentclass[11pt]{article}

\usepackage[margin=1in]{geometry}
\usepackage{amsmath,amssymb,amsthm}
\usepackage{enumitem}
\usepackage{hyperref}
\usepackage{booktabs}
\usepackage{array}
\usepackage{longtable}

\title{\textbf{MATH 617: Number Theory I}\\
Full Course Breakdown}
\author{Johns Hopkins University}
\date{Fall 2023}

\begin{document}

\maketitle

\tableofcontents
\newpage

% ============================================================
\section{Course Overview}
% ============================================================

MATH 617: Number Theory I is a graduate-level course in algebraic
number theory. The course develops the arithmetic of number fields
through the study of Dedekind domains, ideal factorization,
ramification, valuations, local fields, global fields, adèles,
and idèles. If time permits, the course concludes with an
introduction to Tate's thesis and harmonic analysis on local and
global fields.

The conceptual progression of the course is

\[
\boxed{
\text{Number Fields}
\longrightarrow
\text{Dedekind Domains}
\longrightarrow
\text{Ramification}
\longrightarrow
\text{Local Fields}
\longrightarrow
\text{Adèles/Idèles}
\longrightarrow
\text{Tate's Thesis}.
}
\]

% ============================================================
\section{Prerequisites}
% ============================================================

The course assumes graduate-level algebra and analysis,
approximately at the level of qualifying examinations.

\subsection{Algebra}

Students should be familiar with:

\begin{itemize}
    \item Groups, rings, and fields.
    \item Ring homomorphisms and quotient rings.
    \item Ideals and prime/maximal ideals.
    \item Modules and finitely generated modules.
    \item Noetherian rings and modules.
    \item Polynomial rings.
    \item Integral domains and fraction fields.
    \item PID, UFD, and Euclidean domains.
    \item Field extensions.
    \item Algebraic and finite extensions.
    \item Separable and normal extensions.
    \item Galois extensions and Galois groups.
    \item Fundamental theorem of Galois theory.
\end{itemize}

\subsection{Commutative Algebra}

Students should have some familiarity with:

\begin{itemize}
    \item Localization.
    \item Integral extensions.
    \item Integral closure.
    \item Discrete valuation rings.
    \item Fractional ideals.
    \item Prime ideals and maximal ideals.
    \item Modules over Noetherian rings.
    \item Basic dimension theory.
\end{itemize}

\subsection{Analysis}

Students should be familiar with:

\begin{itemize}
    \item Metric spaces and completeness.
    \item Measure theory.
    \item Lebesgue integration.
    \item Complex analysis.
    \item Fourier analysis.
    \item Basic functional analysis.
    \item Topological and locally compact spaces.
\end{itemize}

% ============================================================
\section{Learning Objectives}
% ============================================================

By the end of the course, students should be able to:

\begin{enumerate}
    \item Define and work with number fields.
    \item Determine and use rings of integers.
    \item Compute field norms and traces.
    \item Understand integral extensions and integral closure.
    \item Prove and apply the Dedekind domain properties.
    \item Factor ideals uniquely into prime ideals.
    \item Work with fractional ideals and ideal class groups.
    \item Understand the failure of unique factorization of elements.
    \item Analyze extensions of number fields.
    \item Determine the splitting and ramification of prime ideals.
    \item Work with ramification indices and residue degrees.
    \item Understand differents and discriminants.
    \item Use absolute values and valuations.
    \item Construct and work with completions.
    \item Understand the structure of local fields.
    \item Apply Hensel's lemma.
    \item Distinguish unramified, tame, and wild ramification.
    \item Understand global fields and their places.
    \item Construct the adèle ring.
    \item Construct the idèle group.
    \item Understand the idèle class group.
    \item Relate local and global arithmetic.
    \item If Tate's thesis is covered, apply harmonic analysis
          to local and global fields.
\end{enumerate}

% ============================================================
\section{Module I: Algebraic Number Fields}
% ============================================================

\subsection{Number Fields}

A number field is a finite extension

\[
K/\mathbb{Q}.
\]

Important examples include

\[
\mathbb{Q}(\sqrt{d}),
\qquad
\mathbb{Q}(\sqrt[n]{a}),
\qquad
\mathbb{Q}(\zeta_n).
\]

Topics:

\begin{itemize}
    \item Degree of a number field.
    \item Primitive elements.
    \item Field embeddings.
    \item Real and complex embeddings.
    \item Signature.
    \item Conjugates.
\end{itemize}

\subsection{Integral Elements}

An element $\alpha\in K$ is integral over a ring $R$ if
there exists a monic polynomial

\[
f(x)=x^n+a_{n-1}x^{n-1}+\cdots+a_0
\]

with $a_i\in R$ such that

\[
f(\alpha)=0.
\]

Topics:

\begin{itemize}
    \item Integral elements.
    \item Integral extensions.
    \item Transitivity of integrality.
    \item Integral closure.
\end{itemize}

\subsection{Ring of Integers}

For a number field $K$, define

\[
\mathcal{O}_K
=
\{\alpha\in K:\alpha
\text{ is integral over }\mathbb{Z}\}.
\]

Topics:

\begin{itemize}
    \item Structure of $\mathcal{O}_K$.
    \item Integral bases.
    \item Examples of rings of integers.
    \item Quadratic number fields.
    \item Cyclotomic fields.
\end{itemize}

\subsection{Trace and Norm}

For $K/L$ finite, define

\[
\operatorname{Tr}_{K/L}(\alpha)
=
\sum_{\sigma:K\hookrightarrow \overline{L}}\sigma(\alpha)
\]

and

\[
N_{K/L}(\alpha)
=
\prod_{\sigma:K\hookrightarrow \overline{L}}\sigma(\alpha).
\]

Important properties include

\[
N_{K/L}(\alpha\beta)
=
N_{K/L}(\alpha)N_{K/L}(\beta)
\]

and

\[
\operatorname{Tr}_{K/L}(\alpha+\beta)
=
\operatorname{Tr}_{K/L}(\alpha)
+
\operatorname{Tr}_{K/L}(\beta).
\]

% ============================================================
\section{Module II: Dedekind Domains}
% ============================================================

\subsection{Dedekind Domains}

A Dedekind domain is an integral domain satisfying:

\begin{enumerate}
    \item It is Noetherian.
    \item It is integrally closed.
    \item Every nonzero prime ideal is maximal.
\end{enumerate}

Equivalent local characterization:

\[
R_{\mathfrak p}
\text{ is a discrete valuation ring for every }
\mathfrak p\neq(0).
\]

A fundamental theorem is

\[
\boxed{\mathcal{O}_K \text{ is a Dedekind domain}.}
\]

\subsection{Fractional Ideals}

A fractional ideal of $R$ is an $R$-submodule

\[
I\subset K
\]

such that there exists $a\in R$, $a\neq0$, with

\[
aI\subset R.
\]

Study:

\begin{itemize}
    \item Fractional ideals.
    \item Ideal multiplication.
    \item Invertible ideals.
    \item Principal fractional ideals.
\end{itemize}

\subsection{Unique Factorization of Ideals}

Every nonzero ideal in a Dedekind domain has a unique factorization

\[
I
=
\mathfrak{p}_1^{e_1}
\cdots
\mathfrak{p}_r^{e_r}.
\]

This is one of the central results of the course.

\subsection{Ideal Class Group}

Define

\[
\operatorname{Cl}(K)
=
\frac{\{\text{fractional ideals of }\mathcal{O}_K\}}
{\{\text{principal fractional ideals}\}}.
\]

The class number is

\[
h_K=|\operatorname{Cl}(K)|.
\]

Topics:

\begin{itemize}
    \item Principal ideals.
    \item Ideal classes.
    \item Class groups.
    \item Class numbers.
    \item Failure of unique factorization.
\end{itemize}

% ============================================================
\section{Module III: Geometry of Numbers}
% ============================================================

\subsection{Minkowski Embedding}

Let

\[
r_1=\#\{\text{real embeddings}\},
\qquad
r_2=\#\{\text{pairs of complex embeddings}\}.
\]

Then

\[
[K:\mathbb{Q}]
=
r_1+2r_2.
\]

The Minkowski embedding places $K$ inside

\[
\mathbb{R}^{r_1}\times\mathbb{C}^{r_2}.
\]

\subsection{Lattices}

Study:

\begin{itemize}
    \item Lattices in Euclidean space.
    \item Fundamental domains.
    \item Covolumes.
    \item The lattice associated to $\mathcal{O}_K$.
\end{itemize}

\subsection{Minkowski's Theorem}

Apply Minkowski's theorem to prove finiteness results such as

\[
\boxed{h_K<\infty.}
\]

\subsection{Dirichlet Unit Theorem}

The unit group satisfies

\[
\boxed{
\mathcal{O}_K^\times
\cong
\mu_K\times\mathbb{Z}^{r_1+r_2-1}.
}
\]

Topics:

\begin{itemize}
    \item Roots of unity.
    \item Units.
    \item Rank of the unit group.
    \item Logarithmic embedding.
\end{itemize}

% ============================================================
\section{Module IV: Extensions and Ramification}
% ============================================================

Let

\[
L/K
\]

be a finite extension.

\subsection{Extension of Ideals}

For a prime ideal $\mathfrak{p}\subset\mathcal{O}_K$,

\[
\mathfrak{p}\mathcal{O}_L
=
\mathfrak{P}_1^{e_1}
\cdots
\mathfrak{P}_g^{e_g}.
\]

Define:

\[
e_i=e(\mathfrak{P}_i/\mathfrak{p})
\]

and

\[
f_i
=
[\mathcal{O}_L/\mathfrak{P}_i:
\mathcal{O}_K/\mathfrak{p}].
\]

Under the usual hypotheses,

\[
\boxed{
\sum_{i=1}^g e_i f_i=[L:K].
}
\]

\subsection{Splitting}

Study:

\begin{itemize}
    \item Completely split primes.
    \item Inert primes.
    \item Totally ramified primes.
    \item Partially split primes.
\end{itemize}

\subsection{Ramification}

Study:

\begin{itemize}
    \item Ramification indices.
    \item Residue degrees.
    \item Inertia groups.
    \item Decomposition groups.
    \item Unramified primes.
    \item Ramified primes.
\end{itemize}

% ============================================================
\section{Module V: Discriminants and Differents}
% ============================================================

Given an integral basis
$\alpha_1,\ldots,\alpha_n$, define

\[
\operatorname{disc}
(\alpha_1,\ldots,\alpha_n)
=
\det
\left(
\operatorname{Tr}_{K/\mathbb{Q}}
(\alpha_i\alpha_j)
\right).
\]

The field discriminant is denoted by

\[
d_K.
\]

\subsection{Different}

For an extension $L/K$, study the different ideal

\[
\mathfrak{D}_{L/K}.
\]

Understand the relationships

\[
\boxed{
\text{Different}
\longleftrightarrow
\text{Discriminant}
\longleftrightarrow
\text{Ramification}.
}
\]

% ============================================================
\section{Module VI: Valuations}
% ============================================================

\subsection{Absolute Values}

An absolute value satisfies

\[
|xy|=|x||y|
\]

and

\[
|x+y|\leq |x|+|y|.
\]

Distinguish:

\begin{itemize}
    \item Archimedean absolute values.
    \item Non-archimedean absolute values.
\end{itemize}

\subsection{\texorpdfstring{$p$}{p}-adic Absolute Value}

For

\[
x=p^k\frac{a}{b},
\qquad
p\nmid ab,
\]

define

\[
|x|_p=p^{-k}.
\]

The ultrametric inequality is

\[
|x+y|_p
\leq
\max\{|x|_p,|y|_p\}.
\]

\subsection{Valuations}

A valuation is a map

\[
v:K^\times\rightarrow\mathbb{Z}
\]

satisfying

\[
v(xy)=v(x)+v(y).
\]

The relationship between valuation and absolute value is

\[
|x|_p=p^{-v_p(x)}.
\]

% ============================================================
\section{Module VII: Completions}
% ============================================================

Given a valued field $K$, construct its completion.

Important examples are

\[
\mathbb{Q}
\longrightarrow
\mathbb{R}
\]

and

\[
\mathbb{Q}
\longrightarrow
\mathbb{Q}_p.
\]

Study:

\begin{itemize}
    \item Cauchy sequences.
    \item Complete valued fields.
    \item Extension of valuations.
    \item $\mathbb{Q}_p$.
    \item Finite extensions of $\mathbb{Q}_p$.
\end{itemize}

% ============================================================
\section{Module VIII: Local Fields}
% ============================================================

A local field is a locally compact nondiscrete topological field.

Important examples include

\[
\mathbb{R},\qquad
\mathbb{C},\qquad
\mathbb{Q}_p
\]

and finite extensions of $\mathbb{Q}_p$.

\subsection{Ring of Integers}

For a non-archimedean local field $K$,

\[
\mathcal{O}_K
=
\{x\in K:|x|\leq1\}.
\]

Its maximal ideal is

\[
\mathfrak{p}_K
=
\{x\in K:|x|<1\}.
\]

\subsection{Uniformizer}

Choose $\pi$ such that

\[
\mathfrak{p}_K=(\pi).
\]

Then

\[
x=u\pi^n
\]

for $u\in\mathcal{O}_K^\times$.

Consequently,

\[
K^\times
\cong
\pi^{\mathbb{Z}}
\times
\mathcal{O}_K^\times.
\]

\subsection{Residue Fields}

The residue field is

\[
k_K
=
\mathcal{O}_K/\mathfrak{p}_K.
\]

For finite extensions of $\mathbb{Q}_p$,

\[
|k_K|=p^f.
\]

% ============================================================
\section{Module IX: Extensions of Local Fields}
% ============================================================

For a finite extension $L/K$,

\[
[L:K]=ef.
\]

Study:

\begin{itemize}
    \item Unramified extensions.
    \item Totally ramified extensions.
    \item Tamely ramified extensions.
    \item Wildly ramified extensions.
    \item Extension of valuations.
    \item Local Galois theory.
    \item Ramification groups.
\end{itemize}

\subsection{Hensel's Lemma}

A typical form states that if

\[
f(a)\equiv0\pmod{\mathfrak{p}}
\]

and

\[
f'(a)\not\equiv0\pmod{\mathfrak{p}},
\]

then there exists $\alpha$ satisfying

\[
f(\alpha)=0
\]

and

\[
\alpha\equiv a\pmod{\mathfrak{p}}.
\]

Hensel's lemma is a fundamental tool for lifting solutions
from residue fields to local fields.

% ============================================================
\section{Module X: Global Fields}
% ============================================================

A global field is either:

\[
\text{a number field}
\]

or

\[
\text{a finite extension of }\mathbb{F}_q(t).
\]

Study:

\begin{itemize}
    \item Places.
    \item Valuations.
    \item Completions.
    \item Infinite places.
    \item Finite places.
    \item Local-global relationships.
\end{itemize}

For every place $v$ of $K$, let

\[
K_v
\]

denote the corresponding completion.

% ============================================================
\section{Module XI: Adèles}
% ============================================================

The adèle ring of a global field $K$ is the restricted product

\[
\boxed{
\mathbb{A}_K
=
\prod_v' K_v.
}
\]

An element

\[
x=(x_v)_v
\]

satisfies

\[
x_v\in\mathcal{O}_v
\]

for almost all finite places $v$.

\subsection{Topics}

\begin{itemize}
    \item Restricted products.
    \item Adèle topology.
    \item Diagonal embedding.
    \item The embedding
          \[
          K\hookrightarrow\mathbb{A}_K.
          \]
    \item Quotient $\mathbb{A}_K/K$.
    \item Haar measure.
    \item Characters of the adèle group.
\end{itemize}

% ============================================================
\section{Module XII: Idèles}
% ============================================================

The idèle group is

\[
\boxed{
\mathbb{A}_K^\times.
}
\]

An idèle is a compatible collection

\[
x=(x_v)_v
\]

with

\[
x_v\in K_v^\times
\]

and suitable integrality conditions at almost all finite places.

Study the embedding

\[
K^\times
\hookrightarrow
\mathbb{A}_K^\times
\]

and define the idèle class group

\[
\boxed{
C_K
=
\mathbb{A}_K^\times/K^\times.
}
\]

The idèle class group provides a natural global object
containing information from all local fields simultaneously.

% ============================================================
\section{Module XIII: Tate's Thesis}
% ============================================================

This module is optional and depends on available time.

\subsection{Locally Compact Groups}

Study:

\begin{itemize}
    \item Locally compact groups.
    \item Haar measure.
    \item Characters.
    \item Pontryagin duality.
\end{itemize}

\subsection{Fourier Analysis}

For a locally compact abelian group $G$, define

\[
\widehat{f}(\chi)
=
\int_G
f(x)\overline{\chi(x)}\,dx.
\]

Study:

\begin{itemize}
    \item Fourier inversion.
    \item Fourier transforms.
    \item Poisson summation.
\end{itemize}

\subsection{Schwartz--Bruhat Space}

At archimedean places, use the Schwartz space

\[
\mathcal{S}(\mathbb{R}).
\]

At non-archimedean places, use locally constant,
compactly supported functions.

Globally, these combine into the
Schwartz--Bruhat space

\[
\mathcal{S}(\mathbb{A}_K).
\]

\subsection{Tate's Zeta Integral}

Study integrals of the form

\[
Z(\phi,\chi)
=
\int_{\mathbb{A}_K^\times}
\phi(x)\chi(x)\,d^\times x.
\]

Applications include:

\begin{itemize}
    \item Euler products.
    \item $L$-functions.
    \item Analytic continuation.
    \item Functional equations.
    \item Global-to-local factorization.
\end{itemize}

% ============================================================
\section{Major Conceptual Connections}
% ============================================================

The course should emphasize the following connections.

\subsection{Algebraic and Arithmetic Structure}

\[
\boxed{
\mathcal{O}_K
\rightarrow
\text{Ideals}
\rightarrow
\text{Class Group}.
}
\]

\subsection{Extensions and Ramification}

\[
\boxed{
L/K
\rightarrow
\mathfrak{p}\mathcal{O}_L
\rightarrow
(e,f,g)
\rightarrow
\text{Ramification}.
}
\]

\subsection{Global and Local Number Theory}

\[
\boxed{
K
\rightarrow
K_v
\rightarrow
\text{Local Field}.
}
\]

\subsection{All Local Information Together}

\[
\boxed{
\{K_v\}_v
\rightarrow
\mathbb{A}_K
\rightarrow
\mathbb{A}_K^\times.
}
\]

\subsection{Arithmetic and Analysis}

\[
\boxed{
\mathbb{A}_K
\rightarrow
\text{Harmonic Analysis}
\rightarrow
L\text{-functions}
\rightarrow
\text{Tate's Thesis}.
}
\]

% ============================================================
\section{Difficult Topics}
% ============================================================

The most conceptually difficult parts of the course are expected to be:

\begin{enumerate}
    \item Dedekind domains and ideal factorization.
    \item Fractional ideals and class groups.
    \item Minkowski theory.
    \item Dirichlet's Unit Theorem.
    \item Ramification theory.
    \item Differents and discriminants.
    \item Valuations and completions.
    \item Extensions of local fields.
    \item Hensel's lemma.
    \item Unramified, tame, and wild ramification.
    \item Restricted products and adèles.
    \item Idèles and the idèle class group.
    \item Harmonic analysis on local fields.
    \item Tate's thesis.
\end{enumerate}

% ============================================================
\section{Suggested Weekly Schedule}
% ============================================================

\begin{center}
\begin{tabular}{c l l}
\toprule
\textbf{Weeks} & \textbf{Module} & \textbf{Topics} \\
\midrule
1 & Foundations &
Number fields, embeddings, trace, norm \\
2 & Algebraic integers &
Integrality, rings of integers \\
3 & Dedekind domains &
Ideals, localization, fractional ideals \\
4 & Ideal factorization &
Prime ideals, class groups \\
5 & Geometry of numbers &
Lattices, Minkowski theory \\
6 & Units &
Dirichlet Unit Theorem \\
7 & Extensions &
Extensions of Dedekind domains \\
8 & Ramification &
Splitting, inertia, $e,f,g$ \\
9 & Discriminants &
Different, discriminant, ramification \\
10 & Valuations &
Absolute values, valuations, $p$-adics \\
11 & Completions &
Complete fields, local fields \\
12 & Local fields &
Extensions, residue fields, Hensel's lemma \\
13 & Global fields &
Places and local-global theory \\
14 & Adèles/Idèles &
Restricted products, adèles, idèles \\
15 & Advanced topics &
Tate's thesis / review \\
\bottomrule
\end{tabular}
\end{center}

% ============================================================
\section{Core Theorems to Master}
% ============================================================

A student completing the course should be prepared to understand
and prove, or at least use appropriately, the following major
results:

\begin{enumerate}
    \item Fundamental theorem of Galois theory.
    \item Integrality and transitivity of integrality.
    \item $\mathcal{O}_K$ is a Dedekind domain.
    \item Unique factorization of ideals.
    \item Finiteness of the ideal class group.
    \item Minkowski's theorem.
    \item Dirichlet's Unit Theorem.
    \item Prime decomposition in extensions.
    \item Fundamental relation between $e$, $f$, and $g$.
    \item Properties of discriminants and differents.
    \item Classification of complete non-archimedean fields.
    \item Hensel's lemma.
    \item Structure theory of local fields.
    \item Properties of restricted products.
    \item Basic properties of adèles and idèles.
    \item If covered, Fourier inversion and Poisson summation.
    \item If covered, the basic framework of Tate's thesis.
\end{enumerate}

% ============================================================
\section{Course Dependency Diagram}
% ============================================================

The overall dependency structure can be summarized as

\[
\begin{array}{c}
\text{Graduate Algebra}
\\[4pt]
\downarrow
\\[4pt]
\text{Number Fields}
\\
\downarrow
\\
\text{Algebraic Integers}
\\
\downarrow
\\
\text{Dedekind Domains}
\\
\downarrow
\\
\text{Ideal Factorization}
\\
\downarrow
\\
\text{Extensions of Number Fields}
\\
\downarrow
\\
\text{Ramification}
\\
\downarrow
\\
\text{Valuations}
\\
\downarrow
\\
\text{Completions}
\\
\downarrow
\\
\text{Local Fields}
\\
\downarrow
\\
\text{Global/Local Fields}
\\
\downarrow
\\
\text{Adèles}
\\
\downarrow
\\
\text{Idèles}
\\
\downarrow
\\
\text{Tate's Thesis}
\end{array}
\]

% ============================================================
\section{References}
% ============================================================

\begin{enumerate}
    \item J. Neukirch,
    \emph{Algebraic Number Theory},
    Springer.

    \item A. Weil,
    \emph{Basic Number Theory}.

    \item S. Lang,
    \emph{Algebraic Number Theory},
    Springer.

    \item J. W. S. Cassels and A. Fr\"ohlich,
    \emph{Algebraic Number Theory}.
\end{enumerate}

\end{document}